% move all configuration stuff into one file so we can focus on the content
\documentclass[aspectratio=169,hyperref={pdfpagelabels=false,colorlinks=true,linkcolor=white,urlcolor=blue},t]{beamer}

%%%%%%%%%%%%%%%%%%%%%%%%%%%%%%%%%%%%%%%%%%%%%%%%%%%%%%%%%%%%%%%%%%%%%%%%%%%%%%%%%%
%%%%%%%%%%%%%%%%%%%%%%%%%%%%%%%%%%%%%%%%%%%%%%%%%%%%%%%%%%%%%%%%%%%%%%%%%%%%%%%%%%
% packages
\usepackage{pict2e}
\usepackage{epic}
\usepackage{amsmath,amsfonts,amssymb}
\usepackage{units}
\usepackage{fancybox}
\usepackage[absolute,overlay]{textpos} 
\usepackage{media9} % avi2flv: "C:\Program Files\ffmpeg\bin\ffmpeg.exe" -i TuneFreqFilterbank.avi -b 600k -s 441x324 -r 15 -acodec copy TuneFreqFilterbank.flv
\usepackage{animate}
\usepackage{gensymb}
\usepackage{multirow}
\usepackage{silence}
\usepackage[backend=bibtex,style=ieee]{biblatex}
\AtEveryCitekey{\iffootnote{\tiny}{}}
\addbibresource{references}

%%%%%%%%%%%%%%%%%%%%%%%%%%%%%%%%%%%%%%%%%%%%%%%%%%%%%%%%%%%%%%%%%%%%%%%%%%%%%%%%%%
%%%%%%%%%%%%%%%%%%%%%%%%%%%%%%%%%%%%%%%%%%%%%%%%%%%%%%%%%%%%%%%%%%%%%%%%%%%%%%%%%%
% relative paths
\graphicspath{{graph/}}


%%%%%%%%%%%%%%%%%%%%%%%%%%%%%%%%%%%%%%%%%%%%%%%%%%%%%%%%%%%%%%%%%%%%%%%%%%%%%%%%%%
%%%%%%%%%%%%%%%%%%%%%%%%%%%%%%%%%%%%%%%%%%%%%%%%%%%%%%%%%%%%%%%%%%%%%%%%%%%%%%%%%%
% units
\setlength{\unitlength}{1mm}

%%%%%%%%%%%%%%%%%%%%%%%%%%%%%%%%%%%%%%%%%%%%%%%%%%%%%%%%%%%%%%%%%%%%%%%%%%%%%%%%%%
%%%%%%%%%%%%%%%%%%%%%%%%%%%%%%%%%%%%%%%%%%%%%%%%%%%%%%%%%%%%%%%%%%%%%%%%%%%%%%%%%%
% theme & layout
\usetheme{Frankfurt}
\beamertemplatenavigationsymbolsempty
%\setbeamertemplate{frametitle}[smoothbars theme]
\setbeamertemplate{frametitle}
{
    \begin{beamercolorbox}[ht=1.8em,wd=\paperwidth]{frametitle}
        \vspace{-.1em}%
        \hspace{.2em}{\strut\insertframetitle\strut}
        
        \hspace{.2em}\small\strut\insertframesubtitle\strut
        %\hfill
        %\includegraphics[height=.8cm,keepaspectratio]{CenterMusicTechnology-solid-2lines-white-CoAtag}
        
    \end{beamercolorbox}
    \begin{textblock*}{100mm}(11.6cm,.7cm)
        \includegraphics[height=.8cm,keepaspectratio]{logo_GTCMT_black}
    \end{textblock*}
}

% set this to ensure bulletpoints without subsections
\usepackage{remreset}
\makeatletter
\@removefromreset{subsection}{section}
\makeatother
\setcounter{subsection}{1}

%---------------------------------------------------------------------------------
% appearance
\setbeamercolor{structure}{fg=gtgold}
\setbeamercovered{transparent} %invisible
\setbeamercolor{bibliography entry author}{fg=black}
\setbeamercolor*{bibliography entry title}{fg=black}
\setbeamercolor*{bibliography entry note}{fg=black}

%\usepackage{pgfpages}
%\setbeameroption{show notes}
%\setbeameroption{show notes on second screen=right}
%---------------------------------------------------------------------------------
% fontsize
\let\Tiny=\tiny

%%%%%%%%%%%%%%%%%%%%%%%%%%%%%%%%%%%%%%%%%%%%%%%%%%%%%%%%%%%%%%%%%%%%%%%%%%%%%%%%%%
%%%%%%%%%%%%%%%%%%%%%%%%%%%%%%%%%%%%%%%%%%%%%%%%%%%%%%%%%%%%%%%%%%%%%%%%%%%%%%%%%%
% warnings
\pdfsuppresswarningpagegroup=1
\WarningFilter{biblatex}{Patching footnotes failed}
\WarningFilter{latexfont}{Font shape}
\WarningFilter{latexfont}{Some font shapes}
\WarningFilter{gensymb}{Not defining}



\subtitle{Part 4.1: Instantaneous Features}

%%%%%%%%%%%%%%%%%%%%%%%%%%%%%%%%%%%%%%%%%%%%%%%%%%%%%%%%%%%%%%%%%%%%%%%%%%%%
\begin{document}
    % generate title page
	

\begin{frame}
    \titlepage
    %\vspace{-5mm}
    \begin{flushright}
        \href{http://www.gtcmt.gatech.edu}{\includegraphics[height=.8cm,keepaspectratio]{logo_GTCMT_black}}
    \end{flushright}
\end{frame}


    \section[overview]{lecture overview}
        \begin{frame}{instantaneous features}{overview}
            \begin{itemize}
                \item   \textbf{text book}  
                    \begin{itemize}
                        \item   \href{http://ieeexplore.ieee.org/xpl/ebooks/bookPdfWithBanner.jsp?fileName=6331120.pdf&bkn=6266785&pdfType=chapter}{\underline{\textit{Chapter 3: Instantaneous Features} (pp.~31--63)}}
                    \end{itemize}
                \bigskip
                \item<2->   \textbf{lecture content}
                    \begin{itemize}
                        \item<2->   feature extraction pre-processing
                        \item<3->   statistical features
                        \item<4->   spectral shape features
                        \item<5->   other features
                    \end{itemize}
            \end{itemize}
        \end{frame}

    \section[intro]{introduction}
        \begin{frame}{instantaneous features}{introduction}
            remember the flow chart of a general ACA system:
            \vspace{-2mm}
            \begin{figure}
                \begin{footnotesize}
				\begin{picture}(96,26)
					\setcounter{iXOffset}{0}
					\setcounter{iYOffset}{5}
					\setcounter{iXBlockSize}{28}
					\setcounter{iYBlockSize}{16}
					\setcounter{iYBlockSizeDiv2}{8}
					\setcounter{iDistance}{8}
	
					\addtocounter{iYOffset}{\value{iYBlockSizeDiv2}}
					\addtocounter{iYOffset}{-2}
	
					%\addtocounter{iXOffset}{-1}
					\put(\value{iXOffset}, \value{iYOffset})
						{\text{{\shortstack[c]{audio\\ signal}}}}
					\addtocounter{iXOffset}{1}
	
					\addtocounter{iYOffset}{2}
					\addtocounter{iXOffset}{\value{iDistance}}
	
					\put(\value{iXOffset}, \value{iYOffset})
						{\vector(1,0){\value{iDistance}}}
	
					\addtocounter{iXOffset}{\value{iDistance}}
					\addtocounter{iYOffset}{-\value{iYBlockSizeDiv2}}
					
					\put(\value{iXOffset}, \value{iYOffset})
						{\framebox(\value{iXBlockSize}, \value{iYBlockSize}) {\color{gtgold}{\shortstack[c]{feature\\ extraction}}}}
	
					\addtocounter{iXOffset}{\value{iXBlockSize}}
					\addtocounter{iYOffset}{\value{iYBlockSizeDiv2}}
	
					\put(\value{iXOffset}, \value{iYOffset})
						{\vector(1,0){\value{iDistance}}}
	
					\addtocounter{iXOffset}{\value{iDistance}}
					\addtocounter{iYOffset}{-\value{iYBlockSizeDiv2}}
	
					\put(\value{iXOffset}, \value{iYOffset})
						{\framebox(\value{iXBlockSize}, \value{iYBlockSize}) {\shortstack[c]{decision,\\ interpretation,\\ classification,\\ inference}}}
	
					\addtocounter{iXOffset}{\value{iXBlockSize}}
					\addtocounter{iYOffset}{\value{iYBlockSizeDiv2}}
	
					\put(\value{iXOffset}, \value{iYOffset})
						{\vector(1,0){\value{iDistance}}}
	
					\addtocounter{iXOffset}{\value{iDistance}}
					\addtocounter{iYOffset}{-2}
	
					\addtocounter{iXOffset}{1}
					\put(\value{iXOffset}, \value{iYOffset})
						{\text{{\shortstack[c]{meta\\ data}}}}
					
				\end{picture}
\end{footnotesize}

            \end{figure}
            
            \vspace{-2mm}
            \pause
            \textbf{feature}:
            \begin{itemize}
                \item<2->   \textit{terminology}: 
                    \begin{itemize}
                        \item   audio descriptor
                        \item   instantaneous/short-term/\color<3->{gtgold}{low-level feature}
                    \end{itemize}
                \item<3->   \textit{characteristics}:
                    \begin{itemize}
                        \item	not necessarily musically, perceptually, or semantically meaningful
                        \item	usually one feature per block of data
                    \end{itemize}
            \end{itemize}
        \end{frame}
        \begin{frame}{instantaneous features}{feature}
            \toremember{a feature \ldots}
            
            
            \begin{itemize}
                \item   has to be task-specific,
                \item<2->   is often chosen from a set of established features, but
                \item<3->   can be invented/customized, and
                \item<4->   can also be derived from other data (tags, etc.), or
                \item<5->   from other features.
            \end{itemize}
        \end{frame}
        \begin{frame}{instantaneous features}{feature example}
            waveform envelope of three different signals 
            
            \figwithmatlab{Waveforms}
            
            \vspace{-2mm}
            \begin{columns}
                \column{.16\textwidth}
                \column{.25\textwidth}\centering
                    \includeaudio{excerpt_pop}
                \column{.25\textwidth}\centering
                    \includeaudio{excerpt_stringquartet}
                \column{.25\textwidth}\centering
                    \includeaudio{excerpt_speech}
                \column{.09\textwidth}
            \end{columns}
            
            \bigskip
            \begin{itemize}
                \item<2->   envelopes of waveform seem to have distinct shape
                \item<3->[$\Rightarrow$] a feature describing the envelope shape should help to distinguish these signal types
            \end{itemize}
        \end{frame}

    \section[pre-proc]{audio pre-processing}
        \begin{frame}{instantaneous features}{audio pre-processing}
            \begin{itemize}
                \item   \textbf{goals} (task dependent)
                    \begin{itemize}
                        \item   reduce amount of data (e.g., down-sampling)
                        \item   remove irrelevant information (e.g., surround channels of multi-channel signal)
                        \item   remove information that might impact analysis (e.g., DC offset)
                        \item   increase robustness (e.g., normalization)
                    \end{itemize}
                 %\bigskip
                %\item<2->   \textbf{examples}
                    %\begin{itemize}
                        %\item   downmix to mono
                        %\item   sample rate reduction
                        %\item   filtering (DC removal, HP, BP, ...)
                    %\end{itemize}
            \end{itemize}
        \end{frame}
        \begin{frame}{instantaneous features}{audio pre-processing examples 1/2}
            \begin{itemize}
                \item   \textbf{down-mixing}
                    \[
                        x(i) = \frac{1}{\mathcal{C}}\sum\limits_{c=0}^{\mathcal{C}-1}{x_c(i)} 
                    \]
                    \begin{itemize}
                        \item   \textit{variants}: different channel weights, $\nicefrac{\pi}{2}$ phase shift in one channel, \ldots
                    \end{itemize}
                    
                \bigskip
                \item<2->   \textbf{normalization}
                    \[
                        x(i) = \frac{x_s(i)}{\max\limits_{\forall i}\big(|x_s(i)|\big)} 
                    \]
                    \begin{itemize}
                        \item   \textit{variants}: RMS, LUFS normalization
                        \item<3>   real-time?
                    \end{itemize}
            \end{itemize}
        \end{frame}
        \begin{frame}{instantaneous features}{audio pre-processing examples 2/2}
            \begin{itemize}
                \item   \textbf{filtering}
                    \begin{itemize}
                        \item   \textit{DC removal}
                            \[
                                x(i) = x_{\mathrm{DC}}(i) - \frac{1}{\mathcal{I}}\sum\limits_{i=0}^{\mathcal{I}-1}{x_{\mathrm{DC}}(i)} 
                            \]
                            \begin{itemize}
                                \item   real-time?
                            \end{itemize}
                        \item   other filters: high/band pass, smoothing, ...
                    \end{itemize}
                \bigskip
                \item<2->   \textbf{sample rate reduction}
                \bigskip
                \item<3->   \textbf{quality enhancement} (denoising etc.)
                \bigskip
                \item<3->   \ldots
            \end{itemize}
        \end{frame}

    \section[stats]{statistical features}
        \begin{frame}{instantaneous features}{statistical features}
        \end{frame}

    \section[summary]{lecture summary}

\end{document}

